\documentclass[11pt]{article}
\usepackage[a4paper,margin=22mm]{geometry}
\usepackage[no-math]{fontspec}
\setmainfont{DejaVu Serif}
\usepackage{amsmath,amssymb,amsthm,mathtools,graphicx,xcolor,hyperref,xurl,xspace,ifthen,enumerate}
\setlength{\emergencystretch}{4em}
\SetMathAlphabet{\mathbf}{normal}{OT1}{cmr}{bx}{n}
\SetMathAlphabet{\mathbf}{bold}{OT1}{cmr}{bx}{n}
\newtheorem{theo}{Theorem}[section]
\newtheorem{lemm}[theo]{Lemma}
\newtheorem{thm}[theo]{Theorem}
\newtheorem{proposition}[theo]{Proposition}
\newtheorem{corollary}[theo]{Corollary}
\newtheorem{definition}[theo]{Definition}
\newtheorem{theorem}[theo]{Theorem}
\newtheorem*{remas*}{Remarks}
\newtheorem{remas}[theo]{Remarks}
\newtheorem*{exam*}{Example}
\newtheorem{ques}[theo]{Question}
\providecommand{\og}{«\nobreakspace}
\providecommand{\fg}{\nobreakspace»}
\newtheorem{lem}[theo]{Lemma}
\newtheorem{lemma}[theo]{Lemma}
\newtheorem{prop}[theo]{Proposition}
\newtheorem{coro}[theo]{Corollary}
\newtheorem{corol}[theo]{Corollary}
\newtheorem*{conj*}{Conjecture}
\newtheorem{conj}[theo]{Conjecture}
\newtheorem{Proposition}[theo]{Proposition}
\newtheorem{theoreme}[theo]{Theorem}
\theoremstyle{definition}
\newtheorem{defi}[theo]{Definition}
\newtheorem{exam}[theo]{Example}
\newtheorem{exem}[theo]{Example}
\newtheorem{remark}[theo]{Remark}
\newtheorem{example}[theo]{Example}
\newtheorem{rema}[theo]{Remark}
\newtheorem{nota}[theo]{Notation}
\newtheorem*{rema*}{Remark}
\newtheorem*{defi*}{Definition}
\newtheorem*{theo*}{Theorem}
\newtheorem*{lemm*}{Lemma}
\newtheorem*{prop*}{Proposition}
\newcommand{\enoncename}{}
\newtheorem{corpusenonce}[theo]{\enoncename}
\newenvironment{enonce}[2][]{\renewcommand{\enoncename}{#2}\begin{corpusenonce}}{\end{corpusenonce}}
\newenvironment{enonce*}[2][]{\par\medskip\noindent\textbf{#2.}\itshape}{\par\medskip}
\newenvironment{proofc}[1][\proofname]{\begin{proof}[#1]}{\end{proof}}
\newcommand{\Mc}{\mathcal{M}}
\newcommand{\pointir}{.\ ---\ }
\newcommand{\equalenv}[2]{\expandafter\let\csname#1\expandafter\endcsname\csname#2\endcsname\expandafter\let\csname end#1\expandafter\endcsname\csname end#2\endcsname}
\providecommand{\eg}{e.g.\ }
\providecommand{\ie}{i.e.\ }

\providecommand{\mathof}[1]{\mathrm{#1}}

\numberwithin{equation}{section}


\def\endappendix{}

 \usepackage{mathrsfs}
 \usepackage{tikz-cd}
 \usepackage[poorman]{cleveref}
 \usepackage{pgf,tikz}
\usetikzlibrary{arrows}

\equalenv{Def}{defi}
\equalenv{ex}{exam}
\equalenv{rem}{rema}
\newtheorem{bla}{}
\equalenv{cor}{coro}
\equalenv{con}{conj}
\equalenv{con*}{conj*}
\equalenv{frag}{ques}
\equalenv{lem}{lemm}
\equalenv{lem*}{lemm*}

\Crefname{lemm}{Lemma}{Lemmata}
\Crefname{theo}{Theorem}{Theorems}
\Crefname{defi}{Definition}{Definitions}
\Crefname{exam}{Example}{Examples}
\Crefname{rema}{Remark}{Remarks}
\Crefname{bla}{}{}
\Crefname{coro}{Corollary}{Corollaries}
\Crefname{conj}{Conjecture}{Conjectures}
\Crefname{ques}{Question}{Questions}
\Crefname{nota}{Notation}{Notations}
\Crefname{prop}{Proposition}{Propositions}

\newcommand{\Q}{{\mathbb Q}}
\newcommand{\N}{{\mathbb N}}
\newcommand{\Z}{{\mathbb Z}}
\newcommand{\R}{{\mathbb R}}
\newcommand{\C}{{\mathbb C}}
\newcommand{\HH}{{\mathbb H}}
\newcommand{\cB}{{\mathcal B}}

\newcommand{\supp}{\operatorname{supp}}
\newcommand{\Ima}{\operatorname{Im}}
\newcommand{\rank}{\operatorname{rank}}

\title{Quaternionic unimodular matroids have the half-plane property}
\author{Kummer, Mario, Sawall, David}
\date{doi:10.5802/alco.399; CC BY 4.0}
\begin{document}
\maketitle
\noindent\textit{Editorial note.} Only the designated primary theorem counts. Supporting named results and other statements are uncounted. Focused reading locates and understands the published author proof; it is not a stepwise independent proof verification.\par\medskip
\noindent\textit{Allowed standard prerequisites (editorial).}\begin{itemize}
\item Quaternionic unimodular representation and quaternionic Cauchy-Binet. Formation of M(C), normalized quaternionic matrix with basis minors 0 or1 and quaternionic Cauchy-Binet are prior algebraic inputs; identifying its stable complex determinant with h\_M squared remains the task.
\item Stable determinant polynomials and polynomial identity principle. Determinant of a full-rank complex Gram pencil is stable, and equality on positive orthant implies polynomial equality.
\end{itemize}\medskip
\section{Preliminaries}
We denote by $\N$ and $\N_0$ the set of positive and nonnegative integers, respectively.
In this section we let $E$ always denote a finite set. For $i\in E$ we denote by $\delta_i\in\R^E$ the $i$th unit vector. For $x\in\R^E$ we write $|x|=\sum_{i\in E}|x_i|$. We further let $[n]=\{1,\ldots,n\}$ and denote by $\binom{[n]}{k}$ the set of all $k$-element subsets of $[n]$ for every $k,n\in\N$. We recall the cryptomorphic definitions of M-convex sets and polymatroids.

\begin{Def}
     A subset $J\subseteq\N_0^E$ is \emph{M-convex} if for every $i\in E$ and every $\alpha,\beta\in J$ such that $\alpha_i>\beta_i$, there is $j\in E$ satisfying
            \begin{equation*}
                \alpha_j<\beta_j\textrm{ and }\alpha-\delta_i+\delta_j\in J.
            \end{equation*}
\end{Def}

\begin{Def}
    A \emph{polymatroid} on $E$ is a function $r\colon 2^E\to\N_0$ such that we have for all $S,T\subseteq E$:
            \begin{enumerate}[(i)]
            \item $r(\emptyset)=0$,
                \item $r(S)\leq r(T)$ if $S\subseteq T$, and
                \item $r(S\cup T)+r(S\cap T)\leq r(S)+r(T)$.
            \end{enumerate}
\end{Def}
If $J\subseteq\N_0^E$ is an M-convex set, then we define the function $r_J\colon2^E\to\N_0$ by
    \begin{equation*}
        r_J(S)=\max\{\sum_{i\in S}\alpha_i\mid \alpha\in J\}.
    \end{equation*}
Conversely, if $r\colon 2^E\to\N_0$ is a polymatroid, then we define the set $J_r\subseteq\N_0^E$ by
\begin{equation*}
    J_r=\{x\in\N_0^E\mid \sum_{i\in S}x_i\leq r(S)\textrm{ for all }S\subseteq E\textrm{ and }\sum_{i\in E}x_i=r(E)\}.
\end{equation*}
These two constructions are inverse to each other and define bijections between the set of M-convex sets in $\N_0^E$ and the set of polymatroids on $E$, see for example \cite[\S4.4]{murota}.

\begin{rem}
    A polymatroid $r\colon 2^E\to\N_0$ is the rank function of a matroid $M$ on $E$ if and only if $r(\{i\})\leq1$ for all $i\in E$. In this case we have
    \begin{equation*}
        J_r=\{\sum_{i\in B}\delta_i\mid B\textrm{ is a basis of }M\}.
    \end{equation*}
    We apply the definitions made for polymatroids to matroids by considering their rank functions as polymatroids.
\end{rem}

\subsection{Amalgamation of polymatroids}

\subsection{Stable polynomials}
We briefly recall the different stability and real zero properties of polynomials and their relation to each other. As a general reference we recommend the survey \cite{pemantle}.
\begin{Def}
    Let $P\in\R[x_i\mid i\in E]$ be a polynomial.
    \begin{enumerate}[(a)]
        \item The \emph{support} of $P$ is the unique subset $\supp(P)\subseteq\N_0^E$ such that we can write
        \begin{equation*}
            P=\sum_{\alpha\in\supp(P)}c_\alpha x^\alpha
        \end{equation*}
        for some non-zero $c_\alpha\in\R$.
        \item The polynomial $P$ is called \emph{stable} if for all $z\in\C^E$ such that $\Ima(z_i)>0$ for all $i\in E$ we have $P(z)\neq0$.
        \item The polynomial $P$ is called a \emph{real zero polynomial} if for all $v\in\R^E$ the univariate polynomial $P(t\cdot v)\in\R[t]$ has only real zeros.
        \item Let $P$ be homogeneous and $e\in\R^E$. Then $P$ is called \emph{hyperbolic} with respect to $e$ if for all $v\in\R^E$ the univariate polynomial $P(t\cdot e+v)\in\R[t]$ has only real zeros.
    \end{enumerate}
\end{Def}

Taking $v=0$ in the definition of a real zero polynomial, yields that a real zero polynomial does not vanish at the origin. Similarly, a polynomial $P$ which is hyperbolic in direction $e$, satisfies $P(e)\neq0$. Stability is preserved under taking partial derivatives in coordinate directions, under setting some variables equal to each other and under scaling variables by positive scalars. We summarize the well-known connection between the above concepts.


\begin{Def}
    Let $J\subseteq\N_0^E$ be an M-convex set.
    \begin{enumerate}[(a)]
        \item The \emph{generating polynomial} of $J$ is
        \begin{equation*}
            h_J=\sum_{\alpha\in J}\frac{x^\alpha}{\alpha!}\in \R[x_i\mid i\in E]
        \end{equation*}
        where $\alpha!=\prod_{i\in E}\alpha_i!$.
        \item We say that $J$ has the \emph{half-plane property} if its generating polynomial is stable.
        \item We say that $J$ has the \emph{weak half-plane property} if there exists a homogeneous stable polynomial $P\in\R[x_i\mid i\in E]$ such that $J=\supp(P)$.
        \item We say that a polymatroid $r$ on $E$ has the \emph{(weak) half-plane property} if the associated M-convex set $J_r$ has the (weak) half-plane property.
    \end{enumerate}
\end{Def}

In the remaining part, we observe that in certain nice situations, restricting the polymatroid of a homogeneous stable polynomial corresponds to plugging in zeros for some of the variables.

\section{Quaternionic unimodular matroids}\label{sec:qu}
We recall the definition of quaternionic unimodular (QU) matroids. To this end, let~$\HH$ denote the skew field of quaternions.

\begin{Def}
    Let $E$ be a finite set and $C\subseteq\HH^E$ a submodule of the free left $\HH$-module $\HH^E$. A nonzero element $x\in C$ is called an \emph{elementary chain} of $C$ if $C$ does not contain a nonzero element whose support is strictly contained in the support of~$x$. The submodule $C$ is called \emph{unimodular} if for every elementary chain $x$ of $C$ all nonzero entries of $x$ have the same norm.
\end{Def}

\begin{thm}[{\cite[Theorem~3.7]{skewpartial}}]
    Let $E$ be a finite set and $C\subseteq\HH^E$ unimodular. The set of supports of elementary chains in $C$ is the set of cocircuits of a matroid $M(C)$ on $E$.
\end{thm}

\begin{Def}\label{def:qu}
    A matroid on a finite set $E$ of the form $M(C)$ for some unimodular $C\subseteq\HH^E$ is called \emph{quaternionic unimodular (QU)}.
\end{Def}

The goal of this section is to prove that every QU matroid has the half-plane property. This has been conjectured in \cite[Conjecture~6.9]{skewpartial}.


\begin{Def}[{\cite[page 219]{skewpartial}}]
    Denote by $\varphi\colon\HH\to\C^{2\times 2}$ the map 
\begin{equation*}
    \varphi(a+bi+cj+dk)=
    \begin{pmatrix}
        a+bi& c+di\\
        -c+di& a-bi
    \end{pmatrix}.
\end{equation*}
    We extend $\varphi$ to matrices by applying $\varphi$ entry-wise. Thus for a matrix $A\in\HH^{n\times n}$ we obtain the $2n\times2n$ complex matrix $\varphi(A)$. We define
    \begin{equation*}
      \delta(A)\coloneqq \sqrt{|\det(\varphi(A))|}.  
    \end{equation*}
\end{Def}

\begin{rem}\label{rem:deltaHom}We collect some basic properties of $\delta$.
\begin{enumerate}
    \item For columns $a_1,\ldots,a_n\in\HH^n$ and $\lambda\in\R$, we have
    \begin{equation*}
      \delta(\lambda a_1\ a_2\ldots a_n)=|\lambda|\cdot\delta(a_1\ldots a_n).  
    \end{equation*}
    This follows directly from the definition.
    \item For matrices $A,B\in\HH^{n\times n}$ we have $\delta(AB)=\delta(A)\delta(B)$ and $\delta(A)=\delta(A^*)$. This is \cite[Lemma~5.2]{skewpartial}.
\end{enumerate}
\end{rem}

For a matrix $A\in\HH^{m\times n}$ with $m\leq n$ and $B\subseteq[n]$ of size $m$ we denote by $A[B]$ the $m\times m$ submatrix of $A$ consisting of the columns indexed by $B$.
The following is a version of the Cauchy--Binet theorem over $\HH$.

\begin{prop}[{\cite[Theorem 5.1]{skewpartial}}]\label{prop:CauchyBinetQuat}
    Let $A\in\HH^{m\times n}$ be a matrix over the quaternions and $m\leq n$. Then
    \begin{equation*}
        \delta(AA^*)=\sum_{\substack{B\subseteq[n]\\|B|=m}}\delta(A[B]A[B]^*).
    \end{equation*}
\end{prop}

\begin{lem}\label{lem:basisIffDelta1}
    Let $E$ be a finite set and $C\subseteq\HH^E$ unimodular. Let $d$ be the rank of the matroid $M=M(C)$.
    There is a $d\times |E|$ matrix $A$ over $\HH$ whose rows form a basis of~$C$ such that for all $B\subseteq E$ of size $d$ we have $\delta(A[B])=1$ if $B$ is a basis of $M$ and~$\delta(A[B])=0$ otherwise.
\end{lem}

\begin{proof}
 Let $A$ be a matrix over $\HH$ whose rows form a basis of $C$. By \cite[Lemma~3.14]{skewpartial} it has $d$ rows, so $A\in\HH^{d\times |E|}$. After multiplying $A$ from the left by an invertible $d\times d$ matrix over $\HH$, we can assume by \cite[Corollary~3.26]{skewpartial} that $A$ is a \emph{strong QU matrix} in the sense of \cite[Definition~3.23]{skewpartial}. Now the statement of the lemma follows from {\cite[Claim 5.4.1]{skewpartial}}.
\end{proof}


\begin{thm}[{\cite[Conjecture~6.9]{skewpartial}}]
    Every QU matroid has the half-plane property.
\end{thm}

\begin{proof}
  Let $E=\{1,\ldots,m\}$ and $C\subseteq\HH^m$ be unimodular. Denote by $M=M(C)$ the associated QU matroid and let $A\in\HH^{d\times m}$ be a matrix as in \Cref{lem:basisIffDelta1}. We have to show that the bases generating polynomial $h_M$ of $M$ is stable. This is equivalent to $h_M^2$ being stable. We prove this by showing that $h_M^2$ agrees with the stable polynomial
  \begin{equation*}
    \det\left(\varphi(A)\begin{pmatrix}x_1&&&&\\&x_1&&&\\&&\ddots&&\\&&&x_m&\\&&&&x_m\end{pmatrix}\varphi(A)^*\right).  
  \end{equation*}
    It suffices to show that these two polynomials agree on the positive orthant. We denote by $a_1,\ldots,a_m\in\HH^d$ the columns of $A$. Let $x\in\R_{>0}^m$ and write $x=(x_1^2,\ldots,x_m^2)$ for $x_1,\ldots,x_m\in\R_{>0}$. Then we have
    \allowdisplaybreaks\begin{align}
        h_M(x)^2&=\left(\sum_{\substack{B\subseteq E\\|B|=d}}\delta(A[B]A[B]^*)x^B\right)^2\\
        &=\left(\sum_{\substack{B\subseteq E\\|B|=d}}\delta((x_1a_1\ldots x_ma_m)[B](x_1a_1\ldots x_ma_m)[B]^*)\right)^2\\
        &=\delta((x_1a_1\ldots x_ma_m)(x_1a_1\ldots x_ma_m)^*)^2\\
        &=|\det(\varphi(x_1a_1\ldots x_ma_m)\varphi(x_1a_1\ldots x_ma_m)^*)|\\
        &=|\det\left(\varphi(A)\begin{pmatrix}x_1^2&&&&\\&x_1^2&&&\\
        &&\ddots&&\\&&&x_m^2&\\&&&&x_m^2\end{pmatrix}\varphi(A)^*\right)|\\
        &=\det\left(\varphi(A)\begin{pmatrix}x_1^2&&&&\\&x_1^2&&&\\
        &&\ddots&&\\&&&x_m^2&\\&&&&x_m^2\end{pmatrix}\varphi(A)^*\right).
    \end{align}
    Here we have equality in (2) by \Cref{rem:deltaHom} and \Cref{lem:basisIffDelta1}. (3) holds  by \Cref{rem:deltaHom}. For (4) we use \Cref{prop:CauchyBinetQuat} and (5) follows from the definition of $\delta$. (6) and (7) are obvious.
\end{proof}

\begin{thebibliography}{MMMMMM}
\bibitem[1]{NonPappus} Amini, Nima; Brändén, Petter Non-representable hyperbolic matroids, Adv. Math., Volume 334 (2018), pp. 417-449 \href{https://alco.centre-mersenne.org/articles/10.5802/alco.399/#NonPappus}{Publisher reference}.
\bibitem[2]{BrandenHPP} Brändén, Petter Polynomials with the half-plane property and matroid theory, Adv. Math., Volume 216 (2007) no. 1, pp. 302-320 \href{https://alco.centre-mersenne.org/articles/10.5802/alco.399/#BrandenHPP}{Publisher reference}.
\bibitem[3]{BrandenDLeon} Brändén, Petter; González D’León, Rafael S. On the half-plane property and the Tutte group of a matroid, J. Combin. Theory Ser. B, Volume 100 (2010) no. 5, pp. 485-492 \href{https://alco.centre-mersenne.org/articles/10.5802/alco.399/#BrandenDLeon}{Publisher reference}.
\bibitem[4]{MatroidsArticle} Chen, Justin Matroids: a Macaulay2 package, J. Softw. Algebra Geom., Volume 9 (2019) no. 1, pp. 19-27 \href{https://alco.centre-mersenne.org/articles/10.5802/alco.399/#MatroidsArticle}{Publisher reference}.
\bibitem[5]{HPPlong} Choe, Young-Bin; Oxley, James G.; Sokal, Alan D.; Wagner, David G. Homogeneous multivariate polynomials with the half-plane property, Adv. in Appl. Math., Volume 32 (2004) no. 1-2, pp. 88-187 \href{https://alco.centre-mersenne.org/articles/10.5802/alco.399/#HPPlong}{Publisher reference}.
\bibitem[6]{stickypolymatroids} Csirmaz, László Sticky polymatroids on at most five elements, Studia Sci. Math. Hungar., Volume 58 (2021) no. 1, pp. 136-146 \href{https://alco.centre-mersenne.org/articles/10.5802/alco.399/#stickypolymatroids}{Publisher reference}.
\bibitem[7]{gar} Gȧrding, Lars An inequality for hyperbolic polynomials, J. Math. Mech., Volume 8 (1959), pp. 957-965 \href{https://alco.centre-mersenne.org/articles/10.5802/alco.399/#gar}{Publisher reference}.
\bibitem[8]{masterleon} González D’León, Rafael S. Representing matroids by polynomials with the half-plane property, Master’s thesis, Royal Institute of Technology, Sweden (2009) \href{https://alco.centre-mersenne.org/articles/10.5802/alco.399/#masterleon}{Publisher reference}.
\bibitem[9]{M2} Grayson, Daniel R.; Stillman, Michael E. Macaulay2, a software system for research in algebraic geometry, Available at \url{http://www.math.uiuc.edu/Macaulay2/} \href{https://alco.centre-mersenne.org/articles/10.5802/alco.399/#M2}{Publisher reference}.
\bibitem[10]{ninematroids} Mayhew, Dillon; Royle, Gordon F. Matroids with nine elements, J. Combin. Theory Ser. B, Volume 98 (2008) no. 2, pp. 415-431 \href{https://alco.centre-mersenne.org/articles/10.5802/alco.399/#ninematroids}{Publisher reference}.
\bibitem[11]{murota} Murota, Kazuo Discrete convex analysis, SIAM Monographs on Discrete Mathematics and Applications, Society for Industrial and Applied Mathematics (SIAM), Philadelphia, PA, 2003, xxii+389 pages \href{https://alco.centre-mersenne.org/articles/10.5802/alco.399/#murota}{Publisher reference}.
\bibitem[12]{oxley} Oxley, James Matroid theory, Oxford Graduate Texts in Mathematics, 21, Oxford University Press, Oxford, 2011, xiv+684 pages \href{https://alco.centre-mersenne.org/articles/10.5802/alco.399/#oxley}{Publisher reference}.
\bibitem[13]{pemantle} Pemantle, Robin Hyperbolicity and stable polynomials in combinatorics and probability, Current developments in mathematics, 2011, Int. Press, Somerville, MA, 2012, pp. 57-123 \href{https://alco.centre-mersenne.org/articles/10.5802/alco.399/#pemantle}{Publisher reference}.
\bibitem[14]{skewpartial} Pendavingh, R. A.; van Zwam, S. H. M. Skew partial fields, multilinear representations of matroids, and a matrix tree theorem, Adv. in Appl. Math., Volume 50 (2013) no. 1, pp. 201-227 \href{https://alco.centre-mersenne.org/articles/10.5802/alco.399/#skewpartial}{Publisher reference}.
\bibitem[15]{sticky} Poljak, Svatopluk; Turzíik, Daniel A note on sticky matroids, Discrete Math., Volume 42 (1982) no. 1, pp. 119-123 \href{https://alco.centre-mersenne.org/articles/10.5802/alco.399/#sticky}{Publisher reference}.
\bibitem[16]{purbhoo} Purbhoo, Kevin Total nonnegativity and stable polynomials, Canad. Math. Bull., Volume 61 (2018) no. 4, pp. 836-847 \href{https://alco.centre-mersenne.org/articles/10.5802/alco.399/#purbhoo}{Publisher reference}.
\bibitem[17]{rzamalgamation} Sawall, David; Schweighofer, Markus Amalgamation of real zero polynomials, Indag. Math. (N.S.), Volume 35 (2024) no. 1, pp. 37-59 \href{https://alco.centre-mersenne.org/articles/10.5802/alco.399/#rzamalgamation}{Publisher reference}.
\bibitem[18]{victorsurvey} Vinnikov, Victor LMI representations of convex semialgebraic sets and determinantal representations of algebraic hypersurfaces: past, present, and future, Mathematical methods in systems, optimization, and control (Oper. Theory Adv. Appl.), Volume 222, Birkhäuser/Springer Basel AG, Basel, 2012, pp. 325-349 \href{https://alco.centre-mersenne.org/articles/10.5802/alco.399/#victorsurvey}{Publisher reference}.
\bibitem[19]{wagnerwei} Wagner, David G.; Wei, Yehua A criterion for the half-plane property, Discrete Math., Volume 309 (2009) no. 6, pp. 1385-1390 \href{https://alco.centre-mersenne.org/articles/10.5802/alco.399/#wagnerwei}{Publisher reference}.
\end{thebibliography}
\end{document}
